% Generated draft: requires mathematical and source review.
\documentclass[11pt]{article}
\usepackage[margin=25mm]{geometry}
\usepackage{fontspec}
\setmainfont{Latin Modern Roman}
\usepackage{amsmath,amssymb,mathtools,unicode-math}
\setmathfont{Latin Modern Math}
\usepackage{enumitem,needspace,xurl,hyperref}
\setlength{\parindent}{0pt}
\setlength{\parskip}{5pt}
\setlength{\emergencystretch}{3em}
\newcommand{\N}{\mathbb{N}}
\newcommand{\Z}{\mathbb{Z}}
\newcommand{\Q}{\mathbb{Q}}
\newcommand{\R}{\mathbb{R}}
\newcommand{\C}{\mathbb{C}}
\newcommand{\F}{\mathbb{F}}
\newcommand{\A}{\mathbb{A}}
\newcommand{\PP}{\mathbb{P}}
\newcommand{\E}{\mathbb{E}}
\newcommand{\eps}{\varepsilon}
\newcommand{\dd}{\,\mathrm{d}}
\begin{document}
\section*{UnsolvedMath ID: 30005039}
\textbf{Source problem code:} catalogue entry
\subsection{Definitions and mathematical statement}
Let \(n\geq 2\) be an integer, and let \(F_n\) be the free group of rank \(n\). If \(S\subseteq F_n\) is a generating set, its word metric is
\[
 d_S(x,y)=\min\bigl\{k\in\Z_{\geq 0}:x^{-1}y=s_1\cdots s_k,\quad s_i\in S\cup S^{-1}\bigr\},
\]
where the empty product is the identity. The supplied record does not specify the intended generating set \(S\), restrictions on \(S\), or the category denoted by \(\mathrm{Crs}\).

For clarity, one possible standard metric interpretation can be defined precisely as follows. A map \(f:X\to Y\) between metric spaces is bornologous if
\[
 \forall R\geq 0\ \exists T\geq 0\ \forall x,x'\in X,\qquad
 d_X(x,x')\leq R\ \Longrightarrow\ d_Y(f(x),f(x'))\leq T.
\]
Two maps \(u,v:X\to Y\) are close if \(\sup_{x\in X}d_Y(u(x),v(x))<\infty\). A coarse self-equivalence of \(X\) is a bornologous map \(f:X\to X\) admitting a bornologous map \(g:X\to X\) such that \(g\circ f\) and \(f\circ g\) are close to the identity. Their closeness classes form a group under composition. If \(\mathrm{Crs}\) means this coarse category, then \(\operatorname{Aut}_{\mathrm{Crs}}(F_n,d_S)\) denotes this group for \(X=(F_n,d_S)\). This interpretation is conditional: the supplied excerpt does not establish that it is the intended definition, particularly whether additional group compatibility is required.

\textbf{Requested structural problem.} For \(n\geq 2\), describe the structure of
\[
 \operatorname{Aut}_{\mathrm{Crs}}(F_n,d_S)
\]
and of relevant subgroups. The supplied original excerpt motivates this investigation by asking whether this group, or some subgroup, can be used to study
\[
 \operatorname{Out}(F_n)=\operatorname{Aut}(F_n)/\operatorname{Inn}(F_n),
\]
where \(\operatorname{Inn}(F_n)\) is the subgroup of conjugations by elements of \(F_n\). The motivating analogy is the study of \(\mathrm{GL}(n,\Z)\) through the lattice relationship involving \(\mathrm{SL}(n,\Z)\) and \(\mathrm{SL}(n,\R)\); no corresponding lattice relationship or embedding for \(\operatorname{Out}(F_n)\) is asserted here.

The intended metric, coarse category, subgroups, and meaning of a structural description remain unspecified. Consequently, this is a qualified statement of the supplied research question, not a uniquely determined classification problem. The record labels it open and marks its clean statement as \texttt{corrected\_verified}, while its supplied literature assessment calls the excerpt an incomplete lead-in to Problem 5. These are supplied curation claims, not independent mathematical or source verification.
\subsection{Short English statement}
For free groups \(F_n\) with \(n\geq 2\), investigate the coarse automorphism group and relevant subgroups, motivated by possible applications to \(\operatorname{Out}(F_n)\). The generating set, precise notion of coarse automorphism, and desired structural conclusions require clarification. The supplied open label and source-verification claims have not been independently checked.
\subsection{Sources}
\begin{itemize}
\item Local UnsolvedMath registry: permanent ID \(30005039\), title \emph{Coarse Automorphism Groups of Free Groups}, source problem code \texttt{catalogue entry}, in the supplied Oberwolfach Reports Open Problems record. No individual proposer is supplied.
\item Supplied source citation: \emph{Geometric Structures in Group Theory (2023)}, Oberwolfach Reports, pp.\ 517--576, DOI \texttt{10.4171/owr/2022/11}. \url{https://doi.org/10.4171/owr/2022/11}. This is also the supplied literature-source link labelled ``Oberwolfach report.'' The record separately gives proposed year \(2022\).
\item Supplied literature source labelled ``Coarse invertibility spectrum'': \url{https://doi.org/10.1007/s00209-025-03764-0}. Authors and further bibliographic details are not supplied.
\item Supplied literature assessment dated \texttt{2026-08-22}: the record labels the problem open and describes it as an incomplete lead-in to Problem 5, with the structural classification remaining an active open direction. This is supplied dated curation derived from the OWR material, not an independently verified literature assessment. The record requests independent verification of time-sensitive status and cited claims.
\item The supplied statement-verification field reports that the official report was checked directly and that the clean statement restores definitions and removes extraction artifacts. That verification claim is retained as source metadata; no direct source consultation is claimed in this draft.
\end{itemize}
\end{document}
